% 1 câu đã gỡ khỏi ngân hàng (chuyen_dang) — chép lại vào file dạng để khôi phục

% từ: cau-hoi/lop11/1D7_Dao_ham/1D7-2-6_Gioi_han_ham_so_luong_giac_ham_so_mu_logarit.tex
\begin{ex}%[1D7V2-6]%[ma:3b8f8864]
	Tìm giới hạn $A=\lim\limits_{x  \to+\infty}\left(\dfrac{1+x}{x-1}\right)^x$.
	\loigiai{Với $A=\lim\limits_{x \to+\infty}\left(\dfrac{1+x}{x-1}\right)^x$, đặt $t=x-1$, khi $x \to+\infty$ thì  $t \to+\infty$. Do đó
	$$ 
	A=\lim\limits_{t \to+\infty}\left(\dfrac{t+2}{t}\right)^{t+1}=\lim\limits_{t \to+\infty} \dfrac{t+2}{t}\left(\dfrac{t+2}{t}\right)^t=\lim\limits_{t \to+\infty}  \dfrac{t+2}{t}\cdot \left[\left(1+\dfrac{2}{t}\right)^{\tfrac{t}{2}}\right]^2=\mathrm{e}^2.
	 $$
	 Vậy $A=\mathrm{e}^2$.
	} 
\end{ex}
